\begin{problem}[2019年第36届全国中学生物理竞赛复赛][光纤耦合与微球透镜]{108}
图a是基于全内反射原理制备的折射率阶跃型光纤及其耦合光路示意图。光纤内芯直径 $50\,\mu\mathrm{m}$，折射率 $n_{1}=1.46$；它的玻璃外包层的外径为 $125\,\mu\mathrm{m}$，折射率 $n_{2}=1.45$。氦氖激光器输出一圆柱形平行光束，为了将该激光束有效地耦合进入光纤传输，可以在光纤前端放置一微球透镜进行聚焦和耦合，微球透镜的直径 $D=3.00\,\mathrm{mm}$，折射率 $n=1.50$。已知激光束中心轴通过微球透镜中心，且与光纤对称轴重合。空气折射率 $n_{0}=1.00$。
\begin{subproblem}
为了使光线能在光纤内长距离传输，在光纤端面处光线的最大入射角 $\theta$ 应为多大？
\end{subproblem}
\begin{subproblem}
若光束在透镜聚焦过程中满足近轴条件，为了使平行激光束刚好聚焦于光纤端面（与光纤对称轴垂直）处，微球透镜后表面中心顶点 $O$ 与光纤端面距离应为多大？
\end{subproblem}
\begin{subproblem}
为了使进入光纤的全部光束能在光纤内长距离传输，平行入射激光束的直径最大不能超过多少？
\end{subproblem}
\end{problem}